\documentclass[tikz,border=4pt]{standalone}
\usepackage{liaisons}
\usetikzlibrary{backgrounds}
\begin{document}
\begin{tikzpicture}[x=1cm,y=1cm, solide/.style={line width=1.4pt, line cap=round, line join=round}]
\colorlet{E1}{solideB}\colorlet{E2}{solideA}\colorlet{E3}{solideE}\colorlet{E4}{solideD}
% Symboles (posés d'abord pour définir leurs coordonnées)
\pic[s1=E2, s2=E1, liens=false, name=gA] at (4.5,1.0) {glissiere face};        % E2 coulisse dans E1
\pic[s1=E3, s2=E1, liens=false, name=hB] at (5.0,2.2) {helicoidale face};      % E3 vissée dans E1
\pic[s1=E3, s2=E2, liens=false, name=pC] at (2.3,2.2) {pivot face};            % E3 tourne dans E2
\pic[s1=E4, s2=E3, liens=false, rotate=90, name=dD] at (0.9,2.2) {pivot glissant face}; % E4 coulisse dans la tête de E3, axe y
\pic at (5.8,0.0) {bati};
\pic at (0.5,0.0) {bati};
\begin{scope}[on background layer]
  % E1 : bâti = semelle + mors fixe + support de l'écrou
  \draw[solide, E1] (0.5,0) -- (6.2,0) (3.4,0) -- (3.4,3.2) (4.5,0) -- (gA-s2) (hB-A) ++(0,-0.25) -- ++(0,-0.6) -- (6.2,1.35) -- (6.2,0);
  % E2 : mors mobile = coulisse + bloc vertical
  \draw[solide, E2] (gA-s1b) -- (2.3,1.0) -- (pC-s2) (2.6,1.0) -- (2.6,3.2);
  % E3 : vis, de la tête (D) à l'écrou (B)
  \draw[solide, E3] (dD-s2) -- (hB-s1);
  % E4 : barre de manœuvre, verticale
  \draw[solide, E4] (0.9,0.6) -- (0.9,3.8);
\end{scope}
% Étiquettes
\node[E1, font=\small\bfseries] at (3.75,3.05) {E1};
\node[E2, font=\small\bfseries] at (2.95,3.05) {E2};
\node[E3, font=\small\bfseries, above] at (3.75,2.3) {E3};
\node[E4, font=\small\bfseries, right] at (1.0,3.6) {E4};
\node[font=\small, below] at (4.5,0.7) {A};
\node[font=\small, above] at (5.0,2.5) {B};
\node[font=\small, above] at (2.3,2.5) {C};
\node[font=\small, left] at (0.6,2.2) {D};
\repere{(6.8,0.2)}{x}{y}{1}
\end{tikzpicture}
\end{document}
